\scanpage{42}
\mspage{20}
Reprenons un morphisme généralisé de graphes ordinaires $f:G'\to G^{\deleted{\unsure{'}}}$. On dit que $f$ est un \uline{morphisme de générisation} s'il satisfait la condition 1°) de (48). Notons que $G$ et $f$ sont définis alors à iso près par $G'$ et l'ens $\widetilde B'\subset\widetilde A'$ (stable par $\sigma_{\widetilde A'}$), ou ce qui revient au $\unsure{=|\overline m}$, par $G'$ et la \deleted{\unsure{ens}} partie $B'$ de $A'=\widetilde A'/\sigma_{\widetilde A'}$ ($B'$ est donc un ens. d'arêtes de $G'$). En effet, soit $G'_\bullet$ le sous-graphe de $G'$ défini par $S'$ et $A'\setminus B'\subset A'$ (ou $\widetilde A'\setminus\widetilde B'\subset\widetilde A'$) (i.e. \inserted{par $S'$ et} l'ens des \deleted{arêtes} arcs de $G'$ qui sont au dessus des identités de $\Pi_1(G)$). Alors il résulte de la condition $S'\unsure{\to|{-}}S$ surjectif et les $G'(\alpha)$ connexes ($\Longleftrightarrow$ les $G'(\alpha)$ 0-connexes) que \deleted{les fibres} l'homomorphisme canonique
\[
\pi_0(G'_\bullet)\longrightarrow S
\]
est bijectif, d'où
\[
(49)\qquad S\simeq\pi_0(G'_\bullet),
\]
d'ailleurs on a
\[
(50)\qquad\widetilde A\simeq\widetilde B'\qquad\unsure{\text{compatible}}\text{ avec }\sigma_{\widetilde B'}\text{ et }\sigma_{\widetilde A}
\]
\mspage{41}
